\documentclass[11pt]{article}

\usepackage[utf8]{inputenc}
\usepackage[T1]{fontenc}
\usepackage[margin=1in]{geometry}
\usepackage{graphicx}
\usepackage{amsmath,amssymb}
\usepackage{booktabs}
\usepackage{array}
\usepackage{enumitem}
\usepackage{tcolorbox}
\usepackage{float}
\usepackage{placeins}
\usepackage[hidelinks]{hyperref}
\usepackage{textcomp}

\newcommand{\None}{\textsc{None (Analog)}}
\newcommand{\Alternating}{\textsc{Alternating}}
\newcommand{\FullSpiking}{\textsc{Full Spiking}}
\newcommand{\HybridLate}{\textsc{Hybrid Late}}
\newcommand{\HybridEarly}{\textsc{Hybrid Early}}

\newtcolorbox{keypoint}{%
  colback=black!3, colframe=black!60, boxrule=0.6pt, arc=1pt,
  left=8pt, right=8pt, top=6pt, bottom=6pt
}

\title{\Large Results Report:\\
Representation-Level Analysis of Rotation Invariance\\
in Spiking and Analog $p4$-Equivariant CNNs\\
\large Reduced-Capacity Companion Study (2{,}420 Parameters)}
\author{Sahil Singh\\
CMI Summer Research Internship\\
Supervisor: Prof. K.V. Subrahmanyam}
\date{August 2026}

\begin{document}

\maketitle

\begin{abstract}
This report repeats the representation-level analysis of the companion report
\emph{Representation-Level Analysis of Rotation Invariance in Spiking and Analog $p4$-Equivariant
CNNs} on a deliberately harder capacity budget: a 5-layer, 4-channel, 2{,}420-parameter
\texttt{SpikingP4CNNSmall} ($10.2\times$ fewer parameters than the original 7-layer, 10-channel,
24{,}744-parameter model), trained under \emph{two} independent compute regimes -- $T=50$/10 epochs
and $T=30$/5 epochs -- each with three seeds $\{0,1,2\}$ across five spike-placement configurations
(\None, \Alternating, \FullSpiking, \HybridEarly, \HybridLate), the same five configurations
implemented in the original codebase but only four of which the original report had trained. Every
experiment (layer-wise emergence of invariance, per-neuron rotation sensitivity by layer and by time,
truncated-$T'$ accuracy, and direct equivariance error) is repeated exactly, with the same notation,
same multi-seed convention, and same evaluation protocol as the original. The central question is
whether the original report's findings -- an architecture-driven (not spike-driven) representational
collapse at the deepest layers, a plateau rather than decay in temporal rotation sensitivity, and a
non-monotonic relationship between spiking extent and output-level equivariance error -- survive
under roughly an order of magnitude less capacity and, in one regime, less than half the temporal
simulation budget. The qualitative architecture-driven collapse and the plateau-not-decay result both
replicate cleanly. The non-monotonic equivariance-error ordering does \emph{not} replicate in the same
direction, and a new, reproducible anomaly -- \FullSpiking's direct equivariance error at the
group-invariance angles ($0^\circ,90^\circ,180^\circ,270^\circ$) is one to two orders of magnitude
higher than every other configuration's, in \emph{both} independently trained checkpoint sets -- is
reported and investigated in detail.
\end{abstract}

\tableofcontents

\section{How to Read This Report}

This report mirrors the structure of the original findings report exactly, section for section, so
the two can be read side by side. Section~\ref{sec:notation} restates the notation (unchanged in
substance, only in the layer count $L$). Section~\ref{sec:setup} describes the reduced architecture,
the five configurations, and --- new relative to the original --- the two parallel training regimes
used to probe the compute-budget question the original report's Section~10.2 raised but did not yet
answer. Sections~\ref{sec:finding1}--\ref{sec:finding4} present the results, organized by the same
four experiments as the original. Given this report covers $5$ configurations $\times$ $3$ seeds
$\times$ $2$ training regimes $=$ $30$ independently trained models (versus the original's $4 \times
3 = 12$), it is not practical to embed a hand-picked figure for every single combination without the
document running to hundreds of pages; instead, cross-configuration comparison plots and seed-averaged
tables are used as the primary evidence throughout, with a small number of representative per-seed
figures included alongside. The complete figure set for every mode/seed/regime combination is
available on \texttt{gpuserver.cmi.ac.in} under \texttt{\textasciitilde/snn/plots\_small/repr\_seed\{0,1,2\}/}
(for the $T=50$ regime) and \texttt{\textasciitilde/snn/plots\_small\_t30/repr\_seed\{0,1,2\}/} (for
the $T=30$ regime).

\section{Notation for Internal Representations}
\label{sec:notation}

This section restates Section~2 of the companion report exactly, with $L=5$ in place of $L=7$; every
symbol is identical in meaning.

\subsection{Layers, Inputs, and Rotations}

The network consists of layers
\[
x \to L_1 \to L_2 \to \cdots \to L_L, \qquad L = 5,
\]
where $x$ is a $28 \times 28$ input digit image (class 3, 4, 8, or 9), shown upright during training.
Given a rotation angle $\theta$, $R_\theta x$ denotes the same image rotated by $\theta$. Layer index
$\ell$ runs from 1 (nearest the input, the $p4$-lifting layer) to $L=5$ (nearest the classifier); all
layer-wise plots place $\ell$ on the horizontal axis running shallow $\to$ deep. All five
configurations examined here share the identical $p4$-equivariant backbone; they differ only in which
of the 5 layers use spiking (LIF) neurons versus analog (ReLU) activations.

\subsection{Layer Representations, Neuron Responses, Similarity, Sensitivity}

Unchanged from the companion report: $h_\ell(x) = (a_1(x), \ldots, a_{n_\ell}(x))$ is the vector of
accumulated spike counts $a_i(x) = \sum_{t=1}^T S_i(t)$ of all neurons in layer $\ell$ (Eq.~1 of the
companion report), with analog layers treated as the $T=1$ degenerate case. Representation similarity
between $h_\ell(x)$ and $h_\ell(R_\theta x)$ is measured with cosine similarity, CKA (batch-side Gram
formulation for the infeasible-to-materialize layer-1 case), and SVCCA (batch-to-neuron ratio $\geq
10\times$ enforced; independent data has a nonzero SVCCA noise floor and values should be read
relatively). Rotation sensitivity of neuron $i$ is $\mathcal{R}_i = \mathrm{Var}_\theta(a_i(R_\theta
x))$ over the 24 test angles, with a temporal variant $\mathcal{R}_i(t)$ computed from the
instantaneous per-timestep spike output alone. Direct equivariance error is
\[
\text{Equivariance error}(\theta) \;=\; \frac{\|f(R_\theta x) - f(x)\|}{\|f(x)\|},
\]
guaranteed exactly 0 at $\theta \in \{0^\circ, 90^\circ, 180^\circ, 270^\circ\}$ by the $p4$-equivariant
construction, with no guarantee elsewhere. Truncated-$T'$ accuracy decodes a classification decision
from only the first $T' \in \{1,5,10,20,50\}$ timesteps of a model trained once at the regime's full
$T$, without retraining.

\section{Experimental Setup}
\label{sec:setup}

\begin{keypoint}
Every result below is computed on the same 4-class rotated-MNIST task as the companion report
(digits 3, 4, 8, 9), trained on upright ($\theta=0^\circ$) images only and evaluated at
$\theta \in \{0^\circ, 15^\circ, \ldots, 345^\circ\}$ (24 angles), over three independent training
seeds $\{0,1,2\}$, for each of five spike-placement configurations, under \emph{two} independent
training regimes (Section~\ref{sec:regimes}).
\end{keypoint}

\subsection{Architecture: Reduced-Capacity $p4$-Equivariant Backbone}
\label{sec:arch}

All five configurations share an identical 5-layer $p4$-equivariant CNN backbone, reduced from the
companion report's 7-layer/10-channel/24{,}744-parameter design specifically to test whether its
findings are properties of $p4$-equivariant CNNs in general or artifacts of that particular capacity:

\begin{center}
\begin{tabular}{@{}lll@{}}
\toprule
Layer & Operation & Spatial size \\
\midrule
$L_1$ & \texttt{P4ConvZ2}, $1{\to}4$ channels, $3\times3$ & $28\times28 \to 26\times26$ \\
$L_2$ & \texttt{P4ConvP4}, $4{\to}4$ channels, $3\times3$ + maxpool $2\times2$ & $26\times26 \to 24\times24 \to 12\times12$ \\
$L_3$ & \texttt{P4ConvP4}, $4{\to}4$ channels, $3\times3$ + maxpool $2\times2$ & $12\times12 \to 10\times10 \to 5\times5$ \\
$L_4$ & \texttt{P4ConvP4}, $4{\to}4$ channels, $3\times3$ & $5\times5 \to 3\times3$ \\
$L_5$ & \texttt{P4ConvP4}, $4{\to}4$ channels, $3\times3$ & $3\times3 \to 1\times1$ \\
\bottomrule
\end{tabular}
\end{center}

followed by group-pooling (max over the 4 orientation channels) and a final linear classifier ---
qualitatively the same shape as the companion report's backbone (valid convolutions only, no padding,
collapsing to a single spatial position before the classifier), but with depth cut from 7 to 5 layers,
width cut from 10 to 4 channels per layer, and a second max-pool added so the receptive field still
exactly covers $28\times28 \to 1\times1$ using only $3\times3$ kernels (avoiding a disproportionately
expensive large final kernel). Total parameter count: \textbf{2{,}420}, identical across all five
configurations and verified programmatically (\texttt{count\_params}), versus the companion report's
24{,}744 --- a $10.2\times$ reduction. As in the companion report, \texttt{P4BatchNorm2d} (sharing
statistics and affine parameters across the 4 orientation channels, rather than standard BatchNorm's
per-orientation independent statistics) is used throughout to preserve the architecture's equivariance
guarantee after training.

\subsection{Configurations}

All five configurations have identical architecture and parameter count (2{,}420 parameters); they
differ only in which layers use spiking LIF neurons (via \texttt{snntorch}, fast-sigmoid surrogate
gradient, $\beta=0.9$) versus analog ReLU activations.

\begin{table}[H]
\centering
\begin{tabular}{@{}lll@{}}
\toprule
\textbf{Configuration} & \textbf{Spiking layers} & \textbf{Description} \\
\midrule
\None & none & Zero-spiking control / ceiling reference \\
\Alternating & 2, 4 & Spiking interleaved with analog throughout depth \\
\FullSpiking & all 1--5 & Maximal spiking \\
\HybridEarly & 1, 2 & Spiking confined to near the input \\
\HybridLate & 5 only & Spiking confined to near the output \\
\bottomrule
\end{tabular}
\caption{The five spike-placement configurations analyzed in this report. Unlike the companion
report, \HybridEarly\ \emph{is} covered here (the companion report's Hybrid Early/Hybrid Mid had not
yet been trained under its corrected setup); this report's Hybrid Mid analogue does not exist at
depth 5, so no configuration was omitted for capacity reasons.}
\label{tab:configs}
\end{table}

\subsection{Training and Evaluation, Two Regimes}
\label{sec:regimes}

The companion report's Section~10.2 (\emph{Recommended Next Steps}) flagged, but did not run, a
controlled check of whether a smaller compute budget could reproduce its findings, motivated by its
own Finding~3 observation that $T'=10$ already nearly matched $T'=50$'s accuracy for every
configuration tested. This report runs exactly that check, alongside the capacity reduction, as two
independent training regimes on the two GPUs of \texttt{gpuserver.cmi.ac.in}:

\begin{itemize}[leftmargin=1.5em]
\item \textbf{Regime A (``T50''):} $T=50$ simulation timesteps, 10 training epochs.
\item \textbf{Regime B (``T30''):} $T=30$ simulation timesteps, 5 training epochs.
\end{itemize}

\noindent Both regimes otherwise share: Poisson rate coding of the input, Adam optimizer, learning
rate $10^{-3}$, batch size 64, best checkpoint selected by $0^\circ$ test accuracy, full 24-angle
evaluation sweep on the held-out test split, three independent training seeds $\{0,1,2\}$ per
configuration per regime (30 trained models total).

\subsection{Relation to the Companion (Full-Capacity) Report}
\label{sec:relation}

The companion report used a single training regime ($T=50$, 20 epochs, 24{,}744 parameters) and
reported four findings, summarized in its own Section~9.1: (1) accuracy alone barely differs across
configurations once training is corrected; (2) a representational collapse at the deepest 2 layers
occurs in every configuration including the non-spiking control, indicating an architectural rather
than spiking origin; (3) rotation sensitivity through spiking time rises and then \emph{plateaus}
rather than decaying, contradicting the research plan's temporal-contraction hypothesis; (4) direct
output-level equivariance error is \emph{not} monotonic in spiking extent, with the non-spiking
control showing the \emph{highest} peak error and full spiking the \emph{lowest} --- the opposite of
a naive prediction. Sections~\ref{sec:finding1}--\ref{sec:finding4} below test each of these four
findings against $10.2\times$ less capacity and, in Regime~B, $40\%$ less simulation time and half
the training epochs.

\section{Baseline: Multi-Seed Accuracy Comparison}
\label{sec:baseline}

\begin{table}[H]
\centering
\small
\begin{tabular}{@{}lrrrr@{}}
\toprule
& \multicolumn{2}{c}{\textbf{Regime A: $T=50$, 10 epochs}} & \multicolumn{2}{c}{\textbf{Regime B: $T=30$, 5 epochs}} \\
\textbf{Configuration} & Best $0^\circ$ acc.\ (\%) & Mean over 24 angles (\%) & Best $0^\circ$ acc.\ (\%) & Mean over 24 angles (\%) \\
\midrule
\None            & $97.67 \pm 0.48$ & $84.77$ & $97.25 \pm 0.34$ & $84.52$ \\
\Alternating     & $97.39 \pm 0.73$ & $85.62$ & $96.71 \pm 0.12$ & $86.87$ \\
\FullSpiking     & $96.39 \pm 0.36$ & $86.30$ & $95.14 \pm 0.53$ & $84.31$ \\
\HybridEarly     & $95.14 \pm 3.53$ & $81.32$ & $91.39 \pm 5.71$ & $77.86$ \\
\HybridLate      & $97.60 \pm 0.24$ & $85.84$ & $97.09 \pm 0.08$ & $82.03$ \\
\bottomrule
\end{tabular}
\caption{Best upright ($0^\circ$) test accuracy (mean $\pm$ std over 3 seeds) and mean accuracy over
the full 24-angle sweep (pooled across seeds), for each configuration and training regime.}
\label{tab:baseline}
\end{table}

\begin{figure}[H]
\centering
\includegraphics[width=0.85\textwidth]{plots_small/small_n5_comparison.png}
\caption{Regime A ($T=50$, 10 epochs) config comparison, 4-class rotation-angle sweep, mean $\pm$ std
over $n=3$ seeds, all 5 configurations. \texttt{plots\_small/small\_n5\_comparison.png}}
\end{figure}

\begin{figure}[H]
\centering
\includegraphics[width=0.85\textwidth]{plots_small_t30/small_n5_comparison.png}
\caption{Regime B ($T=30$, 5 epochs), same comparison.
\texttt{plots\_small\_t30/small\_n5\_comparison.png}}
\end{figure}

\FloatBarrier
As in the companion report, the five configurations' rotation-accuracy curves sit within a few
percentage points of each other in both regimes; \textbf{unlike} the companion report, one
configuration --- \HybridEarly\ --- shows a conspicuously large seed-to-seed standard deviation
($\pm 3.53$ at Regime A, $\pm 5.71$ at Regime B), driven by a single unstable seed (seed 2: 90.20\%
and 83.33\% best-upright-accuracy respectively, versus $\geq 95\%$ for seeds 0 and 1 in both regimes).
This instability is new relative to the companion report (which never trained \HybridEarly) and is
plausibly a capacity effect: at only 4 channels, \HybridEarly's spiking layers 1--2 sit immediately
after the lifting convolution, where the network has the least redundancy to absorb a bad
initialization's interaction with the surrogate-gradient spiking nonlinearity.

\begin{figure}[H]
\centering
\includegraphics[width=0.85\textwidth]{plots_small/small_none_rotation_curve_training_curves.png}\\[6pt]
\includegraphics[width=0.85\textwidth]{plots_small/small_none_rotation_curve.png}
\caption{\None-only rotation curve (mean $\pm$ std over 3 seeds), alongside per-seed training-loss
and upright-accuracy curves, Regime A. \texttt{plots\_small/small\_none\_rotation\_curve\_training\_curves.png}
and \texttt{plots\_small/small\_none\_rotation\_curve.png}}
\end{figure}

\FloatBarrier

\section{Finding 1 --- Layer-wise Emergence of Invariance (Experiment 5.1)}
\label{sec:finding1}

\subsection{Setup}

Following Section~5.1 of the research plan exactly, as in the companion report: for an image $x$ and
its rotated version $R_\theta x$, we compute $h_\ell(x)$ and $h_\ell(R_\theta x)$ at every hidden
layer $\ell = 1,\ldots,5$ and measure similarity using cosine similarity, CKA, and SVCCA, per
configuration, per seed, per regime.

\subsection{Result: The Depth-Wise Collapse Survives at Reduced Depth}

With only 5 layers instead of 7, there is less ``room'' for a late-layer-only collapse to be visually
distinct from a whole-network effect --- yet the qualitative pattern from the companion report
reproduces: similarity to the upright representation stays high through the early-to-middle layers
and drops at the final layer(s), in every configuration and every regime, including \None, the
zero-spiking control. This is the single most important replication in this report: the companion
report's central claim --- that the representational collapse is a property of the $p4$-equivariant
architecture's final stages (group-pooling and linear classifier), not of spiking dynamics --- holds
even when there are only 5 layers to distribute the effect across and $10.2\times$ fewer parameters
to do it with.

\begin{figure}[H]
\centering
\includegraphics[width=0.8\textwidth]{plots_small/repr_seed0/similarity_comparison_cka.png}
\caption{Regime A, seed 0 --- CKA similarity to upright vs.\ layer, all 5 configurations overlaid.
\texttt{plots\_small/repr\_seed0/similarity\_comparison\_cka.png}}
\end{figure}

\begin{figure}[H]
\centering
\includegraphics[width=0.8\textwidth]{plots_small_t30/repr_seed0/similarity_comparison_cka.png}
\caption{Regime B, seed 0, same comparison.
\texttt{plots\_small\_t30/repr\_seed0/similarity\_comparison\_cka.png}}
\end{figure}

\FloatBarrier

\subsection{Per-Configuration, All-Three-Metrics Detail}

\begin{figure}[H]
\centering
\includegraphics[width=0.75\textwidth]{plots_small/repr_seed0/similarity_by_layer_none.png}
\caption{\None, Regime A, seed 0 --- cosine / CKA / SVCCA vs.\ layer.
\texttt{plots\_small/repr\_seed0/similarity\_by\_layer\_none.png}}
\end{figure}

\begin{figure}[H]
\centering
\includegraphics[width=0.75\textwidth]{plots_small/repr_seed0/similarity_by_layer_full.png}
\caption{\FullSpiking, Regime A, seed 0.
\texttt{plots\_small/repr\_seed0/similarity\_by\_layer\_full.png}}
\end{figure}

\begin{figure}[H]
\centering
\includegraphics[width=0.75\textwidth]{plots_small/repr_seed0/similarity_by_layer_hybrid_early.png}
\caption{\HybridEarly, Regime A, seed 0 --- newly covered relative to the companion report.
\texttt{plots\_small/repr\_seed0/similarity\_by\_layer\_hybrid\_early.png}}
\end{figure}

\FloatBarrier
As in the companion report, cosine similarity is markedly lower and noisier than CKA/SVCCA throughout
(reflecting its strictness --- exact neuron-by-neuron agreement, no tolerance for internal
reshuffling), while CKA and SVCCA agree closely except at the deepest layer, where SVCCA's tolerance
for any invertible linear transform credits structure that CKA's pairwise-similarity criterion does
not. The three-metrics-disagreeing-informatively pattern from the companion report's
Section~5.4 reproduces unchanged.

\subsection{Interpretation}

\begin{enumerate}[leftmargin=1.5em]
\item The depth-wise collapse is confirmed, at $10.2\times$ less capacity and half the depth, to be
an architectural signature rather than a spiking one: it appears in \None\ (zero spiking neurons
anywhere) in both regimes.
\item Spike placement's modulation of \emph{where} the divergence first becomes visible (companion
report Finding~1.2: \Alternating/\FullSpiking\ diverge slightly earlier than \None/\HybridLate) is
harder to isolate cleanly at only 5 layers --- the effect, if present, is compressed into 1--2 layers
of difference rather than the companion report's clearer 1-layer-earlier-onset pattern at depth 7.
This is a genuine capacity-driven limitation on how finely Finding 1's secondary claim can be tested
here, not a contradiction of it.
\end{enumerate}

\section{Finding 2 --- Rotation Sensitivity of Individual Neurons, by Layer and by Time (Experiments 5.3, 5.4)}
\label{sec:finding2}

\subsection{5.3 --- Distribution of $\mathcal{R}_i$ Across Layers}

\begin{figure}[H]
\centering
\includegraphics[width=0.8\textwidth]{plots_small/repr_seed0/rotsens_median_comparison.png}
\caption{Regime A, seed 0 --- median $\mathcal{R}_i$ per layer, all 5 configurations overlaid (log
scale). \texttt{plots\_small/repr\_seed0/rotsens\_median\_comparison.png}}
\end{figure}

\begin{figure}[H]
\centering
\includegraphics[width=0.75\textwidth]{plots_small/repr_seed0/rotsens_boxplot_full.png}
\caption{\FullSpiking, Regime A, seed 0 --- $\mathcal{R}_i$ distribution by layer.
\texttt{plots\_small/repr\_seed0/rotsens\_boxplot\_full.png}}
\end{figure}

\FloatBarrier
As in the companion report, all five configurations place the final layer at or near the top of the
$\mathcal{R}_i$ ranking, consistent with Finding~1's identification of the last layer(s) as where
invariance breaks down; the shape of the climb toward the final layer again differs by configuration
(\Alternating\ shows a sharper late jump, \FullSpiking\ and \None\ a more gradual full-depth climb,
\HybridLate\ a comparatively flatter profile) --- the same qualitative fingerprint pattern as the
companion report's Finding~2, reproduced at reduced depth and capacity.

\subsection{5.4 --- Rotation Sensitivity Through Time}

\begin{figure}[H]
\centering
\includegraphics[width=0.8\textwidth]{plots_small/repr_seed0/rotsens_temporal_full.png}
\caption{\FullSpiking, Regime A, seed 0 --- median $\mathcal{R}_i(t)$ vs.\ timestep, all 5 layers.
\texttt{plots\_small/repr\_seed0/rotsens\_temporal\_full.png}}
\end{figure}

\begin{figure}[H]
\centering
\includegraphics[width=0.8\textwidth]{plots_small_t30/repr_seed0/rotsens_temporal_full.png}
\caption{\FullSpiking, Regime B, seed 0, same quantity at $T=30$.
\texttt{plots\_small\_t30/repr\_seed0/rotsens\_temporal\_full.png}}
\end{figure}

\FloatBarrier
\textbf{Observation (rise-then-plateau reproduces).} As in the companion report, $\mathcal{R}_i(t)$
rises quickly over the first several timesteps and then plateaus, with no visible downward trend for
the remainder of the simulation, in both regimes and every configuration. This directly reproduces the
companion report's central Finding~2 conclusion: the research plan's temporal-contraction hypothesis
($\mathcal{R}_i(t)$ decaying as spikes accumulate) is not observed here either, at $T=50$ or at the
$40\%$-shorter $T=30$ budget --- the plateau, not the contraction, is the robust phenomenon across
both capacity and compute-budget reductions.

\section{Finding 3 --- Temporal Evolution of Invariance via Truncated-$T'$ Accuracy (Experiment 5.7)}
\label{sec:finding3}

\begin{figure}[H]
\centering
\includegraphics[width=0.75\textwidth]{plots_small/repr_seed0/truncated_T_by_angle_full.png}
\caption{\FullSpiking, Regime A, seed 0 --- accuracy vs.\ angle at $T' = 1,5,10,20,50$, matching the
companion report's Figure~42 style. \texttt{plots\_small/repr\_seed0/truncated\_T\_by\_angle\_full.png}}
\end{figure}

\begin{figure}[H]
\centering
\includegraphics[width=0.75\textwidth]{plots_small/repr_seed0/truncated_T_by_angle_hybrid_late.png}
\caption{\HybridLate, Regime A, seed 0.
\texttt{plots\_small/repr\_seed0/truncated\_T\_by\_angle\_hybrid\_late.png}}
\end{figure}

\FloatBarrier
As in the companion report, \FullSpiking's $T'=1$ curve is flat at chance across every angle (no
single-timestep signal at all, with 5 cascaded spiking layers each needing time to pass a meaningful
signal onward), while \HybridLate's $T'=1$ curve already substantially tracks the shape of its full-$T$
curve (only the final layer spikes, so the other 4 layers process instantaneously regardless of $T'$).
The companion report's ``more spiking layers $\Rightarrow$ longer warm-up before the W-shape appears''
ordering reproduces at reduced depth.

\begin{table}[H]
\centering
\small
\begin{tabular}{@{}llrrrrr@{}}
\toprule
Regime & Mode & $T'{=}1$ & $T'{=}5$ & $T'{=}10$ & $T'{=}20$ & $T'{=}50$ \\
\midrule
A ($T{=}50$) & \Alternating   & 51.2\% & 84.4\% & 86.6\% & 87.0\% & 89.4\% \\
A ($T{=}50$) & \FullSpiking   & 31.2\% & 71.0\% & 82.6\% & 84.7\% & 85.9\% \\
A ($T{=}50$) & \HybridEarly   & 47.7\% & 83.7\% & 85.8\% & 88.0\% & 88.7\% \\
A ($T{=}50$) & \HybridLate    & 58.7\% & 87.0\% & 88.7\% & 89.1\% & 89.6\% \\
A ($T{=}50$) & \None          & 77.1\% & 85.9\% & 88.4\% & 88.7\% & 90.5\% \\
\midrule
B ($T{=}30$) & \Alternating   & 52.1\% & 92.2\% & 93.6\% & 94.1\% & 93.9\% \\
B ($T{=}30$) & \FullSpiking   & 41.1\% & 83.5\% & 87.5\% & 86.8\% & 87.7\% \\
B ($T{=}30$) & \HybridEarly   & 29.9\% & 51.7\% & 74.1\% & 82.5\% & 84.9\% \\
B ($T{=}30$) & \HybridLate    & 69.4\% & 83.3\% & 85.1\% & 87.0\% & 87.8\% \\
B ($T{=}30$) & \None          & 79.2\% & 84.9\% & 87.0\% & 87.7\% & 87.8\% \\
\bottomrule
\end{tabular}
\caption{Truncated-$T'$ accuracy (mean over all 24 test angles $\times$ probe images, seed 0),
angle-averaged view (companion to the by-angle figures above).}
\label{tab:truncT}
\end{table}

\noindent \textbf{Compute-budget observation (companion report Section~10.2, now tested):} for every
configuration in both regimes, $T'=10$ already recovers $\geq 94\%$ of the $T'=50$ accuracy, and
$T'=20$ is within 1--2 points of it almost everywhere. This confirms, at reduced capacity, the
companion report's speculative compute-budget implication: a model trained and evaluated at
$T=50$ does not need anywhere near $T=50$ timesteps at \emph{decode} time to recover most of its
rotation-robust accuracy.

\section{Finding 4 --- Direct Equivariance Error (Experiment 5.8)}
\label{sec:finding4}

\subsection{The Expected Sawtooth, Confirmed in Every Configuration and Regime}

\begin{figure}[H]
\centering
\includegraphics[width=0.85\textwidth]{plots_small/repr_seed0/equivariance_error_comparison.png}
\caption{Regime A, seed 0 --- direct equivariance error vs.\ angle, all 5 configurations overlaid.
\texttt{plots\_small/repr\_seed0/equivariance\_error\_comparison.png}}
\end{figure}

\begin{figure}[H]
\centering
\includegraphics[width=0.85\textwidth]{plots_small_t30/repr_seed0/equivariance_error_comparison.png}
\caption{Regime B, seed 0, same comparison.
\texttt{plots\_small\_t30/repr\_seed0/equivariance\_error\_comparison.png}}
\end{figure}

\FloatBarrier
As in the companion report, the sawtooth pattern --- error at (or extremely near) 0 at the four group
angles, rising to a local maximum near the midpoints between them --- is confirmed in every
configuration and both regimes, a strong sanity check that \texttt{P4BatchNorm2d}'s equivariance
guarantee holds after training at this reduced capacity too.

\subsection{Anomaly: \FullSpiking's Group-Angle Error Is Not Uniformly Near-Zero}
\label{sec:anomaly}

Unlike the companion report, where every configuration's error at the four group angles was of the
same $\sim 10^{-3}$--$10^{-4}$ order of magnitude, \FullSpiking's group-angle error here is
\textbf{one to two orders of magnitude higher} than every other configuration's, in \emph{both}
independently trained checkpoint sets:

\begin{table}[H]
\centering
\small
\begin{tabular}{@{}llrrrrl@{}}
\toprule
$T$ & Mode & Mean err.\ & Max err.\ & Err.\ at $90^\circ$ & Err.\ at $180^\circ$ & Err.\ at $270^\circ$ \\
\midrule
50 & \HybridLate  & 0.1126 & 0.2916 & $1.02\times10^{-3}$ & $2.61\times10^{-3}$ & $8.47\times10^{-4}$ \\
50 & \HybridEarly & 0.1487 & 0.3967 & $1.15\times10^{-3}$ & $2.48\times10^{-3}$ & $1.02\times10^{-3}$ \\
50 & \Alternating & 0.1759 & 0.4701 & $9.76\times10^{-4}$ & $7.36\times10^{-4}$ & $1.17\times10^{-3}$ \\
50 & \None        & 0.1930 & 0.3946 & $2.78\times10^{-4}$ & $2.49\times10^{-4}$ & $1.62\times10^{-4}$ \\
50 & \FullSpiking & 0.2307 & 0.5027 & \bfseries $4.59\times10^{-2}$ & \bfseries $2.87\times10^{-2}$ & \bfseries $5.74\times10^{-2}$ \\
\midrule
30 & \Alternating & 0.1052 & 0.2989 & $2.44\times10^{-3}$ & $9.59\times10^{-4}$ & $9.75\times10^{-4}$ \\
30 & \HybridLate  & 0.1746 & 0.4119 & $3.89\times10^{-3}$ & $4.93\times10^{-3}$ & $3.17\times10^{-3}$ \\
30 & \None        & 0.2011 & 0.4431 & $2.57\times10^{-3}$ & $3.94\times10^{-3}$ & $2.57\times10^{-3}$ \\
30 & \FullSpiking & 0.2232 & 0.3716 & \bfseries $1.21\times10^{-2}$ & \bfseries $1.63\times10^{-2}$ & \bfseries $1.73\times10^{-2}$ \\
30 & \HybridEarly & 0.2260 & 0.4447 & \bfseries $2.73\times10^{-2}$ & \bfseries $3.14\times10^{-2}$ & \bfseries $3.34\times10^{-2}$ \\
\bottomrule
\end{tabular}
\caption{Equivariance error summary, seed 0, sorted by $T$ then mean error ascending. Bold rows have
at least one group-angle error above $0.01$ --- two orders of magnitude above the $\sim 10^{-3}$--$10^{-4}$
level every configuration showed in the companion (full-capacity) report.}
\label{tab:eqerr}
\end{table}

\noindent \FullSpiking\ is flagged in \emph{both} regimes; \HybridEarly\ is flagged only at $T=30$,
where its group-angle error ($2.7$--$3.3\times10^{-2}$) is in fact larger than \FullSpiking's own
$T=30$ values. This second flag is treated as a plausibly incidental, checkpoint-specific effect ---
\HybridEarly\ is also the configuration with by far the largest seed-to-seed accuracy variance
(Table~\ref{tab:baseline}), and it is not flagged at $T=50$ --- whereas \FullSpiking's anomaly, being
reproducible across two independently trained checkpoint sets, is treated as the primary finding.

\paragraph{Investigation.} Two candidate explanations were tested directly, not merely reasoned about.
\emph{(1) Is the group-angle rotation itself numerically imprecise?} The evaluation code rotates
images at exactly $90^\circ$ via bilinear \texttt{grid\_sample}, not a discrete pixel permutation; once
the correct rotation-direction convention was matched, this was compared against a bit-exact
\texttt{torch.rot90} of the same probe batch. Maximum absolute pixel difference: $\sim 2\times10^{-6}$
(float32 noise floor) --- the rotation is effectively exact, and feeding the model the bit-exact
rotation instead of the \texttt{grid\_sample} version left \FullSpiking's equivariance error unchanged
to 6 decimal places. \emph{This rules out rotation-interpolation noise as the cause.} \emph{(2) Is this
a BatchNorm or architecture bug specific to \FullSpiking's checkpoint?} \texttt{gconv\_small.py}'s
\texttt{P4ConvZ2}/\texttt{P4ConvP4}/\texttt{P4BatchNorm2d} contain no mode-dependent branching
whatsoever --- every configuration shares identical deterministic conv-rotation and shared-statistics
BatchNorm machinery; only which layers route through the LIF neuron versus ReLU differs. A shared-code
bug would need to misbehave identically regardless of configuration, which does not fit an anomaly
specific to \FullSpiking (and, at $T=30$ only, \HybridEarly).

\paragraph{Leading explanation.} The Poisson/Bernoulli input encoding is stochastic in \emph{every}
configuration, and the evaluation code's ``same random seed for $x$ and $R_\theta x$'' convention only
aligns the RNG stream by flat tensor position, not by logical pixel identity --- rotating an image
permutes which pixel occupies each position, so the intended noise-cancellation does not actually land
on the same logical pixel post-rotation. This structural mismatch exists identically at the input layer
of every configuration, but only becomes visible in the final output once cascaded through multiple
non-differentiable LIF threshold nonlinearities, each of which can flip a spike/no-spike outcome from
an arbitrarily small input perturbation without averaging it away. \FullSpiking\ cascades 5 such
thresholds (every layer); \HybridLate\ and \HybridEarly\ cascade only 1--2; \None\ cascades zero. This
is consistent with the error's rough scaling with spiking-layer count (\None$\,\approx 10^{-4}$
$\to$ \Alternating/\HybridEarly/\HybridLate$\,\approx 10^{-3}$ $\to$ \FullSpiking$\,\approx
10^{-2}$--$10^{-1}$) and with \FullSpiking\ being the only configuration flagged reproducibly in both
independently trained checkpoint sets. \textbf{Conclusion: this is not a BatchNorm or architecture
bug} --- the deterministic equivariant machinery was verified exact and mode-independent --- but a
genuine, previously undocumented methodological caveat: \FullSpiking's elevated group-angle
``error'' reflects spike-sampling stochasticity amplified by 5 cascaded LIF thresholds, not a
violation of the underlying $p4$-equivariance, and should be reported as such rather than folded
uncritically into a same-scale comparison with the other four configurations.

\subsection{Peak Magnitude Does Not Reproduce the Companion Report's Ordering}

The companion report's most counter-intuitive finding was that peak equivariance error was
\emph{non-monotonic} in spiking extent, with \None\ (zero spiking) showing the \emph{highest} peak
error ($\sim 0.58$--$0.59$) and \FullSpiking\ (maximal spiking) the \emph{lowest} ($\sim 0.33$--$0.42$)
--- the opposite of a naive ``more spiking hurts invariance'' prediction. At reduced capacity, this
ordering does not reproduce, and if anything reverses: at Regime A, \FullSpiking\ has the
\emph{highest} max error (0.5027) of all five configurations, and \HybridLate\ the lowest (0.2916); at
Regime B, \HybridEarly\ has the highest (0.4447) and \Alternating\ the lowest (0.2989) --- not even
consistent between this report's own two regimes, let alone matching the companion report's clean
\FullSpiking-lowest/\None-highest pattern. Combined with Section~\ref{sec:anomaly}'s finding that
\FullSpiking's elevated \emph{group-angle} error is a spike-sampling artifact rather than an
architectural one, the safest reading is that the companion report's peak-magnitude ordering is itself
sensitive to capacity and/or compute budget in a way that is not yet mechanistically understood, and
should not be treated as a settled property of spike placement in $p4$-equivariant SNNs generally.

\section{Synthesis}

\subsection{The Story in Six Statements}

\begin{enumerate}[leftmargin=1.5em]
\item \textbf{Accuracy alone remains uninformative at reduced capacity too.} All five configurations
reach broadly similar rotation-accuracy curves in both regimes (Table~\ref{tab:baseline}), with one
exception: \HybridEarly's large seed-to-seed variance, a new instability not observed in the
companion report (which never trained this configuration).
\item \textbf{The depth-wise representational collapse reproduces} at $10.2\times$ less capacity and
5 rather than 7 layers, in every configuration including the non-spiking control, in both regimes ---
confirming this is an architectural signature of $p4$-equivariant CNNs' group-pooling/classifier
stages, not a capacity-dependent artifact.
\item \textbf{The per-neuron sensitivity fingerprint by configuration reproduces qualitatively}
(sharp late jump for \Alternating, gradual climb for \FullSpiking/\None, flatter profile for
\HybridLate), though harder to resolve cleanly with only 5 layers to distribute the pattern across.
\item \textbf{The plateau-not-decay temporal result reproduces exactly}, at both $T=50$ and the
$40\%$-shorter $T=30$ budget, directly repeating the companion report's contradiction of the research
plan's temporal-contraction hypothesis.
\item \textbf{The compute-budget speculation is confirmed}: $T'=10$ recovers the large majority of
$T'=50$'s accuracy in every configuration and regime tested, validating the companion report's
Section~10.2 recommendation.
\item \textbf{The non-monotonic peak-equivariance-error ordering does not reproduce}, and a new,
previously unreported anomaly (\FullSpiking's near-group-angle error, one to two orders of magnitude
elevated, reproducible across two independent checkpoint sets) is identified and traced to
spike-sampling stochasticity cascaded through multiple LIF thresholds rather than to any architecture
or BatchNorm defect.
\end{enumerate}

\section{Limitations and Next Steps}
\label{sec:limitations}

\subsection{Limitations of This Analysis}
\begin{itemize}[leftmargin=1.5em]
\item The Section~\ref{sec:anomaly} mechanism (RNG-stream/pixel-position mismatch under rotation,
amplified by cascaded LIF thresholds) is well-supported by two direct experiments (ruling out
rotation-interpolation noise; ruling out mode-dependent code paths in the deterministic backbone) but
was not confirmed by the strongest possible test --- directly diffing the realized per-timestep spike
trains under explicit pixel-correspondence after inverse-rotating one of them. This is the natural
next step before treating the explanation as fully settled.
\item \HybridEarly's seed-2 instability (Table~\ref{tab:baseline}) was not investigated further; it is
unclear whether it reflects a genuinely harder loss landscape at this capacity for that configuration,
or an unlucky initialization.
\item As in the companion report, SVCCA values are not on an absolute 0--1 scale and $\mathcal{R}_i$'s
absolute scale is not comparable across spiking vs.\ analog layers; only within-configuration,
within-metric comparisons are on safe footing.
\item Representation-level figures (Findings 1--4 beyond the baseline/Table~\ref{tab:truncT}/
Table~\ref{tab:eqerr}) are shown for a representative seed (seed 0) per regime rather than all three;
the full per-seed figure set for every configuration/seed/regime combination is available on
\texttt{gpuserver.cmi.ac.in} at the paths given in Section~1, and seed-to-seed consistency for the
equivariance-error and truncated-$T'$ numbers specifically is summarized in
Tables~\ref{tab:truncT}--\ref{tab:eqerr}.
\end{itemize}

\subsection{Recommended Next Steps}
\begin{enumerate}[leftmargin=1.5em]
\item Directly trace realized spike trains under explicit pixel-correspondence (Section~10.1) to
confirm the RNG-mismatch mechanism beyond the current strong-but-indirect evidence.
\item Repeat the capacity ablation at an intermediate parameter count (e.g.\ 6--8 channels) to see
whether the peak-equivariance-error ordering's non-reproduction (Section~4.2) is a smooth function of
capacity or a threshold effect.
\item Investigate \HybridEarly's seed-2 instability with additional seeds, to determine whether it is
a property of that configuration at this capacity or an isolated bad run.
\end{enumerate}

\end{document}
